\section{Comparison of operational models}
\label{sec:comparison96}\label{sec:comparison84}
The feasible set, not merely the word ``memory'', determines a decision
problem.  We record four comparisons relevant to the general and
spectral results above.

\paragraph{Constrained-separability testers.}
Ohst, Zhang, Nguyen, Pl\'avala and Quintino
\cite[Sections~5.3--5.4 and~6.1--6.2]{OZNPQ89} formulate auxiliary
memory restrictions by positive separable decompositions with affine
conditions on their factors.  Their Definition~23 and Theorem~24
characterize complete measure-and-reprepare adaptive testers and their
constrained-separability formulation.  Our class $\mathfrak C$ uses
that complete-readout boundary.  A reset tester retaining an old
receiver instead need not have separable final complete-call effects:
the later joint readout is unrestricted, while coherent import into
the new source is prohibited.  Our fixed-initial objective constrains
the common Gram barycenter and evaluates it explicitly for the
shared-basis payoff.  The general tester framework is an antecedent,
not a claim superseded by a first construction of memory hierarchies.

\paragraph{Autonomous coherent memory.}
Zonnios and Binder \cite[Proposition~2, Theorem~1 and Corollary~1]{ZB90}
optimize over a repeated stationary instrument carrying a coherent
register between steps and retaining the classical record.  Their
finite-time completeness compiles a general adaptive tester into an
autonomous machine with sufficiently large memory and a counter.
Here the source cannot receive the old coherent register, public
controls may depend on the step, and the two specified receiver cuts
are not an all-time workspace dimension.  Saturation at
$\operatorname{rank}\rho$ in our biased experiment concerns the
exact value of that family with its first acquisition fixed.  It is
not their arbitrary-horizon compilation theorem.

\paragraph{Majorization and ensemble transformations.}
Nielsen \cite{Nielsen94} characterizes deterministic pure-state
transformations by majorization.  Jonathan and Plenio
\cite[Theorem~1]{JonathanPlenio94} treat ensembles of pure-state
outputs through averaged Schmidt constraints.  These are direct
antecedents for the spectral barycenter step.  Our feasible
instruments act only on the already prepared receiver, have the common
sum $\sum_hA_{yh}=\rho$ after each first device label, and optimize
subsequent measurement decisions.  Their transformation targets are
not the payoff matrices \eqref{eq:gameblocks91}.  The least-majorant
argument, finite convexity and support inverse are used here with
credit; the specific contribution is their combination with the exact
rank-constrained filtered-swap leaf and its message comparator.

\paragraph{Assistance and concave roofs.}
Entanglement assistance \cite{Assistance94} and general roof methods
\cite{Uhlmann91} optimize averages over decompositions of a fixed
state.  The present all-density and pure-atom roofs have such a
mathematical form, but their atoms, rank cuts and continuation reward
come from the reset normal form.  The uniform message criterion asks
whether these continuation functions are concave; the binary witness
changes one receiver-to-source classical permission, not an
entanglement-conversion task.  A universal majorant is a condition on
all admissible atoms and cannot be replaced by finite sampling.

Finally, a measurement channel with classical output is not a L\"uders
instrument exposing a post-measurement state \cite{EQ90}.
The local geometry in the supplement, the general variational theory,
and the closed shared-basis formulas remain separate theorem families.
No priority for majorization, concave envelopes, swap tests, or
memory-constrained discrimination as general mechanisms is claimed.
